%%
%% EPITECH PROJECT, 2025
%% Tek2 Internship Report
%% File description:
%% Mission chapter of the
%% report
%%


\section{Mission Générale}

\begin{flushleft}
  La mission générale du stage, telle qu'elle m'a été
  confiée de TrackIt, était de réaliser des projets
  pour notre client, étant principalement du Back-End
  en Python utilisant la plateforme Cloud AWS, et
  toujours lié à la gestion des médias.
  Ainsi, je devais non-seulement faire preuve de
  compétences techniques, mais aussi de transparence
  et de professionalisme envers le client, comme nous
  étions en contacte direct pendant la durée du stage.
\end{flushleft}

\begin{flushleft}
  En pratique, le client me chargeait d'une nouvelle
  tâche, voire d'un nouveau projet, à chaque fois que
  je finissais un projet et qu'il avait été approuvé
  par le reste de l'équipe. Ainsi, j'ai pu réaliser
  cinq grands projets lors de la durée de mon stage,
  qui vont être expliqués plus en détail par la suite.
\end{flushleft}

\begin{figure}[h]
  \centering
  \includegraphics[width=8cm]{project_flow}
  \caption{Déroulement simplifié d'un projet}
\end{figure}

\begin{flushleft}
  Le diagramme si-dessus montre les étapes simplifiés du
  déroulement d'un projet avec le client. Concrètement,
  on implémentait le projet (seul ou en groupe selon le
  projet), puis on testait le projet pour s'assurer que
  tout fonctionne conformément au cahier des charges.
  Ensuite, avec l'équipe entière, nous avions approuvé
  le projet pour assurer que la qualité du projet soit à
  la hauteur d'un projet de TrackIt. Finalement, les deux
  ingénieurs de TrackIt, Timothé CONIEL et Mathieu BRUJAN,
  ont déployé le projet en production chez notre client,
  pour qu'il puisse s'en servir.
\end{flushleft}

\newpage

\section{Venues Workflow}

\begin{flushleft}
  En début de stage, mon collègue Omar AHMED et moi avions
  été chargés de réaliser un projet ensemble. Ce projet,
  intitulé \textit{Venues Workflow} (anglais: flux de travail
  de sites d'évènements), consistait à automatiser une suite
  de tâches que notre client devait faire manuellement.
\end{flushleft}

\begin{flushleft}
  Comme dit dans la première partie, le client organise des
  évènements à grande échelle. Ainsi, une équipe du client
  est chargée de photographier les lieux de l'évènement,
  afin de pouvoir tout planifier en fonction de la place
  disponsible. Or, l'équipe chargée de prendre ces
  photographies n'ont pas les accès à la l'entièreté de la
  plateforme MAM, ce qui faisait que JS devait manuellement
  récupérer les photographies, les extraire, télécharger les
  données dans l'emplacement voulu (de façon à ce qu'ils soient
  accessibles dans le MAM), puis rajouter les métadonnées et la
  vignette par soi-même.
\end{flushleft}

\begin{flushleft}
  Ainsi, le client nous demandait de créer un \textit{workflow}
  (\sccad{} une suite de tâches qui se lancent automatiquement
  lorsqu'une action spécifique est faite), qui faisait les tâches
  à sa place (automatiquement) lorsque un membre de l'équipe des
  photographes télécharge sa collection de photographies dans un
  dossier spécifique.
\end{flushleft}

\begin{figure}[h]
  \centering
  \includegraphics[width=12cm]{venues_workflow}
  \caption{Déroulement du Workflow}
\end{figure}

\begin{flushleft}
  Comme le montre le diagramme ci-dessus, le workflow débute
  quand un membre de l'équipe de photographie télécharge une
  fichier zip dans un dossier dédiée aux sites d'évènements.
  Une fois téléchargé, une fonction hébergée sur AWS est lancée,
  qui est responsable d'extraire le fichier zip et d'afficher
  le progrès du traitement en parallèle, afin que l'utilisateur
  sache que le traitement est en cours. Une fois terminé,
  une nouvelle ressource liée aux données extraites est crée dans
  la plateforme MAM, contenant des métadonnées automatiquement
  assignées (comme le nom du lieu, la date, le pays, \scetc).
  Finalement, une image vignette est automatiquement rajoutée à
  la ressource, et le workflow est terminé.
\end{flushleft}

% \section{Private PIP}

\section{MAM Shortcuts Extension}

\begin{flushleft}
  Le deuxième projet du stage, nommé \textit{MAM Shortcuts Extension}
  (anglais: extension de raccourcis MAM), est une extension de navigateur
  permettant d'attribuer des raccourcis de clavier à des actions réalisables
  sur le site de la plateforme MAM.\@
\end{flushleft}

\begin{flushleft}
  En effet, le site de la plateforme MAM ne supporte pas les raccourcis clavier
  typiques (\sccad{} les flèches pour naviguer, Ctrl + C pour copier, Ctrl + V
  pour coller, \scetc). Ainsi, l'équipe du client, étant responsable de naviguer
  la plateforme et manipuler les ressources médias qui y sont présent, se retrouvait
  à faire des mouvements de souris lent et répétitifs, alors qu'ils pourraient être
  simplifiés avec des raccourcis clavier.
\end{flushleft}

\begin{flushleft}
  Ainsi, on m'a confié le projet de réaliser cette extension
  de navigateur (spécifique aux navigateurs Chromium,
  comme Google Chrome et Bing, entre autre).
\end{flushleft}

\begin{figure}[h]
  \centering
  \includegraphics[width=6cm]{extension_home}
  \caption{Page d'accueil de l'extension}
\end{figure}

\begin{center}
  \textit{Les parties floues sont liées à la plateforme MAM du
    client et ne sont donc pas divulgués.}
\end{center}

\begin{flushleft}
  L'image ci-dessus montre les quatre pages de l'extension: la
  page \textit{raccourcis}, contenant les raccourcis clavier;
  la page \textit{actions}, \scou{} les actions qui s'activent
  lorsqu'un raccourci clavier est actionné sont définies; la
  page \textit{paramètres}, sur laquelle on peut configurer
  l'extension; et finalement la page \textit{aide} qui explique
  comment on peut utiliser l'extension, ainsi que ce qu'elle
  permet de faire.
\end{flushleft}

\begin{flushleft}
  En considérant que les utilisateurs qui naviguent la plateforme
  pourraient vouloir des raccourcis pour des actions auxquelles
  nous n'avions pas encore pensés, j'ai choisi d'implémenter un
  système de création d'action qui permet aux utilisateurs
  de définir leurs propres actions, voire de personnaliser
  les actions qui existent déjà, à leur propre guise.
\end{flushleft}

\begin{figure}[h]
  \centering
  \includegraphics[width=5cm]{extension_create_action}
  \caption{Menu de création d'action de l'extension}
\end{figure}

\begin{flushleft}
  Comme le montre l'image ci-dessus, le créateur d'actions permet
  de définir ses propres actions en définissant le type d'action,
  les paramètres de l'action, et en configurant d'autres options.
\end{flushleft}

\begin{flushleft}
  Une fois qu'une action est crée, elle peut être exportée sous
  forme d'un fichier (comme le montre l'image ci-dessous), ce qui
  permet aux utilisateurs de se partager des actions
  entre\textendash{}eux, ce qui est pratique pour ceux qui ne
  souhaitent pas prendre le temps d'en créer.
\end{flushleft}

\begin{figure}[h]
  \centering
  \includegraphics[width=5cm]{extension_action_menu}
  \caption{Page d'actions de l'extension}
\end{figure}

\newpage

\section{MAM SDK}

\section{MAM MCP Server}
